\documentclass[11pt,a4paper]{article}
\usepackage[utf8]{inputenc}
\usepackage[T1]{fontenc}
\usepackage[margin=2.5cm]{geometry}
\usepackage{amsmath,amssymb}
\usepackage{booktabs}
\usepackage{longtable}
\usepackage{microtype}
\usepackage{graphicx}
\usepackage{listings}
\usepackage{color}
\usepackage{hyperref}
\usepackage{xcolor}
\usepackage{siunitx}

\definecolor{codegray}{rgb}{0.95,0.95,0.95}
\definecolor{codeblue}{rgb}{0.1,0.2,0.6}
\definecolor{darkgray}{rgb}{0.35,0.35,0.35}
\definecolor{commentgreen}{rgb}{0.0,0.45,0.0}
\definecolor{warnorange}{rgb}{0.75,0.35,0.0}

\lstset{
    backgroundcolor=\color{codegray},
    basicstyle=\small\ttfamily,
    breaklines=true,
    columns=fullflexible,
    commentstyle=\color{commentgreen}\itshape,
    frame=single,
    keepspaces=true,
    keywordstyle=\color{codeblue}\bfseries,
    language=bash,
    showstringspaces=false,
    numbers=none,
    xleftmargin=6pt
}

\hypersetup{
    colorlinks=true,
    linkcolor=codeblue,
    urlcolor=codeblue,
    pdftitle={Damping Sweep \& Analysis Campaign Manual --- compass\_sim},
    pdfauthor={J. P. Sinnecker / Antigravity IDE}
}

\newcommand{\code}[1]{\texttt{#1}}
\newcommand{\file}[1]{\texttt{#1}}
\newcommand{\param}[1]{\texttt{--#1}}
\newcommand{\warn}[1]{\textcolor{warnorange}{\textbf{#1}}}

\title{\textbf{Damping Sweep \& Analysis Campaign Manual}\\[1ex]
\large\textit{compass\_sim --- scripts \file{damping\_sweep.py},
\file{damping\_sweepV02.py},\\
\file{sweep\_damping\_hysteresis.py},
\file{damping\_sweep\_analysis.py}}\\[1ex]
\normalsize Manual revision: July 2026}
\author{J.~P.~Sinnecker}
\date{\today}

\begin{document}

\maketitle
\tableofcontents
\newpage

\begin{quote}
\warn{Note (added after the fact):} of the four scripts this manual covers,
only \file{damping\_sweep\_analysis.py} is unaffected by the following.
\file{damping\_sweep.py} and \file{damping\_sweepV02.py} call functions
(\code{make\_grid()}, \code{relax()}, \code{label\_magnetic\_domains()})
that no longer exist in the engine and cannot run at all (see
\file{docs/AUDIT.md} bug B1). \file{sweep\_damping\_hysteresis.py} is also
non-functional outside \param{dry\_run} (bug B3). None of this manual
covers \file{damping\_sweepV03.py}, written after this manual and the only
campaign wrapper actually compatible with the current \file{compass.py} --
see \file{USER\_GUIDE.md} \S9 for it instead. \file{damping\_sweep\_analysis.py}
itself assumes the pre-V03 CSV "S" column convention (a sparse event count)
and silently misinterprets V03's \code{flip\_field} column if pointed at
V03 output (bug B2) -- it has not been updated for V03's schema. The
content below is left as a historical record of the V01/V02 design, not a
current usage guide.
\end{quote}

% ======================================================================
\section{Overview}
\label{sec:overview}
% ======================================================================

This manual describes the Python simulation-campaign suite built to study
the central scientific question of the \texttt{compass\_sim} project:

\begin{quote}
\itshape ``Does the viscous damping coefficient control the universality
class of avalanches and the shape of the hysteresis loop in classical
two-dimensional dipolar arrays, and does this depend on the lattice
geometry (square, triangular, honeycomb)?''
\end{quote}

\noindent The suite consists of four scripts:

\begin{center}
\begin{tabular}{lp{10cm}}
\toprule
Script & Role \\
\midrule
\file{damping\_sweep.py} & Sequential campaign runner (hysteresis + relaxation) \\
\file{damping\_sweepV02.py} & Parallel campaign runner (\texttt{ProcessPoolExecutor});
  same physics as V01, adds \param{n\_workers} and \param{resume} \\
\file{sweep\_damping\_hysteresis.py} & Earlier standalone script; now physically
  consistent with V01/V02 after corrections \\
\file{damping\_sweep\_analysis.py} & Post-processing: metrics, power-law fitting,
  waiting-time distribution, autocorrelation, plots \\
\bottomrule
\end{tabular}
\end{center}

\medskip
\textbf{The sweep execution proceeds in two stages per configuration:}
\begin{description}
\item[Stage 1 --- Hysteresis.] A five-ramp cycle
  ($0 \to +B_{\max} \to 0 \to -B_{\max} \to 0 \to +B_{\max}$)
  sweeps the external field to trace the hysteresis loop and capture
  avalanche events.
\item[Stage 2 --- Free relaxation.] The external field is set to zero and the
  system relaxes freely from its post-hysteresis state; magnetic-domain sizes
  and counts are measured in the final configuration.
\end{description}

% ======================================================================
\section{Physical Model and Parameter Derivation}
\label{sec:physics}
% ======================================================================

\subsection{Equations of motion}

Each needle $i$ satisfies Newton's second law for rotation:
\begin{equation}
  I\,\ddot{\theta}_i = \tau_i^{\mathrm{dip}} + \tau_i^{\mathrm{ext}} - b\,\dot{\theta}_i
\end{equation}
where $I$ is the moment of inertia [kg\,m$^2$], $b$ is the viscous damping
coefficient [N\,m\,s/rad], and the torques arise from dipolar coupling and
the external field.

\subsection{Quality factor $Q$}

The dimensionless quality factor characterises the dynamical regime:
\begin{equation}
  Q = \frac{\omega_0 I}{b}, \qquad
  \omega_0 = \sqrt{\frac{m B_{\mathrm{ref}}}{I}}, \qquad
  B_{\mathrm{ref}} = \frac{\mu_0}{4\pi}\frac{2m}{r_{\mathrm{nn}}^3}
\end{equation}
where $m$ is the needle magnetic moment [A\,m$^2$], $r_{\mathrm{nn}}$ is the
nearest-neighbour distance ($= 2R$ for all implemented geometries), and
$B_{\mathrm{ref}}$ is the dipolar field between nearest neighbours.

\medskip
\noindent Three dynamical regimes are identified:

\begin{center}
\begin{tabular}{lll}
\toprule
Regime & Condition & Physical behaviour \\
\midrule
Underdamped & $Q \gg 1$ & Needles oscillate before settling; chain-reaction avalanches \\
Critically damped & $Q \approx 1$ & Fastest settling with no overshoot \\
Overdamped & $Q \ll 1$ & Slow monotonic relaxation; narrow hysteresis loops \\
\bottomrule
\end{tabular}
\end{center}

\subsection{Needle geometry and physical parameters}

All scripts derive moment of inertia $I$ and magnetic moment $m$ from the
physical geometry of the needle blade and its pivot pin using
\texttt{compass.py}'s authoritative functions:

\begin{align}
  &\text{\code{compute\_inertia\_from\_geometry}(needle\_len, needle\_width, thickness,
    density, pivot\_*)} \label{eq:inertia}\\
  &\text{\code{compute\_moment\_from\_geometry}(needle\_len, needle\_width, thickness, Ms)}
    \label{eq:moment}
\end{align}

\noindent where the needle geometry is derived from the circumscribed radius $R$
and the needle fraction $f$:
\begin{equation}
  L_{\mathrm{needle}} = f \times 2R, \qquad
  W_{\mathrm{needle}} = 0.22 \times L_{\mathrm{needle}}, \qquad
  f = 0.80\ \text{(default)}
\end{equation}

\noindent The pivot pin contributes an additive term to $I$ (Eq.\,\ref{eq:inertia});
if \param{pivot\_radius}~$= 0$ (the default), the pivot contribution is zero.

\warn{Important}: in earlier versions of \file{sweep\_damping\_hysteresis.py},
the module-level fallback constants \code{cs.MOMENT\_DEFAULT} and
\code{cs.INERTIA\_DEFAULT} were used instead of the geometry-derived values.
This is now fixed --- all scripts call the same geometry functions as
\code{compass.py\,main()}.

% ======================================================================
\section{\texttt{damping\_sweep.py} --- Serial Campaign}
\label{sec:v01}
% ======================================================================

\subsection{Operation}

\file{damping\_sweep.py} imports \texttt{compass.py} and calls its
\code{relax()} function for each combination:
\[
  \text{geometry} \times \text{damping/}Q \times \text{seed}
\]
Runs are executed sequentially on a single CPU core.

The damping grid is built once using the \textit{square} geometry as reference
and then the same physical damping values $b$ are applied to the other
geometries.  Because $r_{\mathrm{nn}}$ (and therefore $Q$) differs slightly
between geometries, the actual $Q$ of every run is computed and stored in its
metadata JSON.

\subsection{Command-line interface}

\begin{longtable}{p{4.6cm}p{2.0cm}p{8.0cm}}
\toprule
Flag & Default & Description \\
\midrule
\endhead
\param{out\_dir} & \textit{required} & Output directory \\
\param{grid\_n} & 30 & Grid dimension $N = M$ (total needles $= N^2$) \\
\param{n\_seeds} & 5 & Random seeds per (geometry, damping) \\
\param{n\_dampings} & 8 & Points in the $Q$ grid \\
\param{Q\_min} & 0.05 & Minimum quality factor \\
\param{Q\_max} & 15.0 & Maximum quality factor \\
\param{geometries} & \textit{all three} & Comma-separated: \texttt{square,triangular,honeycomb} \\
\param{R} & 0.025 & Circumscribed needle radius [m] \\
\param{B\_max\_factor} & 8.0 & $B_{\max} = \text{factor} \times B_{\mathrm{ref}}$ \\
\param{t\_sim\_periods} & 40.0 & Simulation duration in units of $T_0 = 2\pi/\omega_0$ \\
\param{pbc} & 0 & Periodic boundary conditions (0 = off) \\
\param{use\_gpu} & 0 & GPU acceleration via CuPy (0 = off) \\
\param{skip\_relax\_stage} & 0 & Omit Stage~2 (domain statistics) \\
\param{damping\_noise} & 0.0 & Relative per-needle damping variation (e.g.\ 0.05 = 5\%) \\
\param{needle\_frac} & 0.80 & Blade length fraction of $2R$ \\
\param{needle\_thickness} & 0.0004 & Steel blade thickness [m] \\
\param{steel\_density} & 7850.0 & Blade steel density [kg\,m$^{-3}$] \\
\param{steel\_Bsat} & \textit{None} & Override saturation flux density [T] \\
\param{moment} & \textit{None} & Override magnetic moment [A\,m$^2$] \\
\param{inertia} & \textit{None} & Override moment of inertia [kg\,m$^2$] \\
\param{pivot\_radius} & 0.0 & Pivot pin radius [m] \\
\param{pivot\_thickness} & 0.0 & Pivot pin thickness [m] \\
\param{pivot\_density} & 8500.0 & Pivot pin material density [kg\,m$^{-3}$] \\
\param{pivot\_mass} & \textit{None} & Override pivot mass directly [kg] \\
\param{dt\_factor} & 0.05 & Time step $\Delta t$ as fraction of $T_0$ \\
\param{noise} & 1.5 & Initial angle noise [rad] \\
\param{domain\_tol} & 15.0 & Angular tolerance for domain clustering [\textdegree] \\
\param{quick\_test} & 0 & If 1: $N=12$, 1 seed, 3 dampings, 8 periods --- pipeline check \\
\bottomrule
\end{longtable}

\subsection{Output structure}

\begin{lstlisting}
<out_dir>/
+-- manifest.csv              # one row per run (stage, tag, Q, paths, ...)
+-- hysteresis/
|   +-- square_damp00_seed00.csv   # columns: t, B, M, S
|   +-- ...
+-- relax/
|   +-- square_damp00_seed00.csv
|   +-- ...
+-- meta/
    +-- square_damp00_seed00_hyst.json   # full physical parameters + Q
    +-- square_damp00_seed00_relax.json  # + domain_sizes list
\end{lstlisting}

CSV column names are \texttt{t, B, M, S} where \texttt{S} is the integer
spin-flip count at each integration step (number of needles that reversed
orientation).

% ======================================================================
\section{\texttt{damping\_sweepV02.py} --- Parallel Campaign}
\label{sec:v02}
% ======================================================================

V02 retains the full parameter set of V01 and adds two features:

\subsection{Parallel execution with worker sandboxing}

\texttt{ProcessPoolExecutor} distributes runs across \param{n\_workers}
processes.  Each worker is initialised in a private scratch directory
(\texttt{\_scratch/worker\_xxxx/}) so that \texttt{hysteresis\_loop.csv}
and other simulator side-effect files do not collide.

\warn{Never combine \texttt{--n\_workers $>$ 1} with \texttt{--use\_gpu 1}}:
multiple processes cannot safely share a single GPU context.

\subsection{Incremental manifest and resume}

\texttt{manifest.csv} is appended and flushed after \textit{each} completed
run.  An interrupted campaign can be resumed with \param{resume~1}; the script
reads the existing manifest, identifies completed (stage, tag) pairs, and
skips them.

\subsection{V02-specific flags}

\begin{tabular}{p{3.5cm}p{2.2cm}p{9cm}}
\toprule
Flag & Default & Description \\
\midrule
\param{n\_workers} & 1 & Number of parallel processes \\
\param{resume} & 0 & If 1, skips runs already in \texttt{manifest.csv} \\
\bottomrule
\end{tabular}

\subsection{Suppression of \texttt{compass.py} terminal output}

Each call to \code{relax()} inside V01 and V02 is wrapped with
\code{contextlib.redirect\_stdout(open(os.devnull,~`w'))}, so that
\texttt{compass.py}'s own status lines
(\textit{``Pre calculated Dipolar Tensor...''},
\textit{``[Ctrl+I to stop...]''},
\textit{``Hysteresis data saved: hysteresis\_loop.csv''})
are discarded during campaign runs.
This is essential for the progress bar to work correctly: any
\texttt{\textbackslash n} printed to stdout advances the terminal cursor and
breaks the \texttt{\textbackslash r} overwrite trick used by the bar.
Stderr output (Python exceptions, warnings) is \textit{not} redirected and
reaches the terminal normally; errors are therefore still visible.

% ======================================================================
\section{Terminal Progress Bar}
\label{sec:progressbar}
% ======================================================================

All three campaign/analysis scripts share an identical helper function
\code{\_bar(done, total, label, extra, width=36)} that renders a
one-line overwriting progress bar:

\begin{lstlisting}
  HYSTERESIS [########################............]    2/3   66.7%  square_damp01_seed00  Q=0.866  0.3s
  RELAX      [####################################]    3/3  100.0%  square_damp02_seed00  Q=0.050  nd=105
  CAMPAIGN   [##############################......]    5/6   83.3%  HYST square_damp02_seed00  Q=0.050  [1 active]
\end{lstlisting}

\noindent (In the actual terminal the filled blocks are rendered as \texttt{U+2588} and the
empty blocks as \texttt{U+2591}; in the parallel bar suffix the hourglass \texttt{U+231B}
precedes the active-worker count.)

\noindent The bar uses Unicode block characters (U+2588 filled / U+2591 light shade);
no external dependencies are required.  Behaviour:
\begin{itemize}
  \item Called with \texttt{done=0} before the loop to show the empty bar;
    called after each completed job with the updated count.
  \item While $\texttt{done} < \texttt{total}$, the line is overwritten
    via \texttt{\textbackslash r} (carriage return) --- the cursor returns
    to column~0 and the new bar overwrites the previous one on the same
    terminal line.
  \item When $\texttt{done} = \texttt{total}$ the line ends with
    \texttt{\textbackslash n}, advancing to the next line.
\end{itemize}

\noindent The \texttt{extra} field shows context specific to each script:

\begin{center}
\begin{tabular}{lll}
\toprule
Script & Label & Extra info \\
\midrule
\file{damping\_sweep.py} hysteresis & \texttt{HYSTERESIS} & last tag $\cdot$ $Q$ $\cdot$ wall time \\
\file{damping\_sweep.py} relaxation & \texttt{RELAX} & last tag $\cdot$ $Q$ $\cdot$ \texttt{n\_domains} \\
\file{damping\_sweepV02.py} serial & \texttt{CAMPAIGN} & stage $\cdot$ tag $\cdot$ $Q$ $\cdot$ wall time \\
\file{damping\_sweepV02.py} parallel & \texttt{CAMPAIGN} & stage $\cdot$ tag $\cdot$ $Q$ $\cdot$ \texttt{[N active]} \\
\file{damping\_sweep\_analysis.py} hyst & \texttt{HYSTERESIS} & current tag \\
\file{damping\_sweep\_analysis.py} relax & \texttt{RELAX} & current tag \\
\bottomrule
\end{tabular}
\end{center}

\warn{Note}: the progress bar requires that \texttt{compass.py}'s own stdout
be suppressed (see Section~\ref{sec:v02}).  If compass.py prints
a newline between two bar updates the \texttt{\textbackslash r} trick fails
and each update appears on a new line instead of overwriting the previous one.

% ======================================================================
\section{\texttt{sweep\_damping\_hysteresis.py} --- Standalone Script}
\label{sec:standalone}
% ======================================================================

This is an earlier standalone script that runs only Stage~1 (hysteresis) with
hard-coded geometry and a per-run scratch directory.  Its output is compatible
with \file{damping\_sweep\_analysis.py}.

\subsection{Corrections applied (July 2026)}

Three bugs present in the original version were corrected:

\begin{enumerate}
  \item \textbf{CUTOFF fixed}: the interaction cutoff was \texttt{8.0} (treated
    as metres by \code{relax()}, which uses units of $r_{\mathrm{nn}}$, giving
    a physically incorrect $8\,r_{\mathrm{nn}}$ cutoff instead of the intended
    $3.5\,r_{\mathrm{nn}}$).  Corrected to \texttt{CUTOFF = 3.5}, consistent
    with V01/V02 and the \texttt{compass.py} CLI default.

  \item \textbf{MOMENT/INERTIA fixed}: the script used
    \code{cs.MOMENT\_DEFAULT} and \code{cs.INERTIA\_DEFAULT} (module-level
    fallback values) instead of calling the geometry derivation functions.
    Now uses \code{cs.compute\_moment\_from\_geometry()} and
    \code{cs.compute\_inertia\_from\_geometry()}, making results
    cross-comparable with V01/V02.

  \item \textbf{$B_{\max}$ and CSV columns fixed}: \code{B\_MAX} was a hard-coded
    constant; it is now computed dynamically as \texttt{B\_MAX\_FACTOR~$\times$~$B_{\mathrm{ref}}$}.
    CSV column names corrected from \texttt{B\_proj, M\_proj} to \texttt{B, M}
    so that \file{damping\_sweep\_analysis.py} can read the files without
    modification.  A \texttt{damp\_idx} field was also added to metadata.
\end{enumerate}

\subsection{Limitations}

\file{sweep\_damping\_hysteresis.py} runs sequentially (no \param{n\_workers})
and does not generate Stage~2 relaxation data.  For new campaigns, prefer
\file{damping\_sweepV02.py}.

% ======================================================================
\section{\texttt{damping\_sweep\_analysis.py} --- Post-Processing}
\label{sec:analysis}
% ======================================================================

The analysis script reads the directory produced by any of the sweep scripts
(identified via \texttt{manifest.csv}) and computes all metrics.

\subsection{Hysteresis loop metrics}

\begin{description}
  \item[Loop area.]  $\displaystyle A = \tfrac{1}{2}\left|\sum_{k} B_k\,M_{k+1} - B_{k+1}\,M_k\right|$
    (shoelace formula applied to the closed loop in the $B$--$M$ plane).
    Measures energy dissipated per cycle.

  \item[Coercive field $B_c$.]  Linear interpolation of $B$ where $M$ crosses
    zero.  The median over all such crossings in the cycle is reported.

  \item[Remanence $M_r$.]  Linear interpolation of $M$ at the time-indices
    where $B = 0$.  The median is reported.
\end{description}

\subsection{Avalanche detection --- two methods}

Both methods are applied to every hysteresis run; their results are saved
separately so that the two approaches can be compared.

\paragraph{Method 1: $dM/dB$ proxy (MAD-based).}
The field cycle is split into monotonic ramps at its turning points.  Within
each ramp, the derivative $dM/dB$ is computed and a robust threshold is set
using the Median Absolute Deviation:
\begin{equation}
  \text{threshold} = \operatorname{median}(dM/dB) + \kappa \cdot \operatorname{MAD}(dM/dB)
\end{equation}
where $\kappa = $ \param{mad\_threshold} (default 4.0).  Contiguous steps above
the threshold form one avalanche; its size is $s = \sum |dM|$ over those steps.

\paragraph{Method 2: spin-flip counter (\texttt{S} column).}
A cleaner alternative using the integer \texttt{S} column written by
\texttt{compass.py} at each integration step.  An avalanche is a maximal run
of consecutive steps with $S > 0$; its size is $\sum S$ over the run.
This method involves no numerical differentiation and is immune to numerical
drift near inflection points.

The S-column method also returns the \textbf{inter-arrival times}
(waiting times between consecutive avalanche starts), used in
Section~\ref{sec:waiting_time}.

\subsection{Power-law exponent $\tau$ --- MLE fit (Clauset--Shalizi--Newman)}

To stabilise the MLE estimate, avalanche sizes are \textit{aggregated across
all seeds} for each (geometry, damp\_idx) combination before fitting.  The
procedure follows Clauset, Shalizi \& Newman 2009:

\begin{enumerate}
  \item For candidate lower bounds $x_{\min}$ (quantiles of the sample), compute
    the MLE estimator for the continuous power-law exponent:
    \begin{equation}
      \hat{\tau} = 1 + n\left[\sum_{i=1}^n \ln\!\left(\frac{s_i}{x_{\min}}\right)\right]^{-1}
    \end{equation}
  \item Evaluate the KS distance $D$ between the empirical CDF and the
    theoretical power-law CDF.
  \item Select the $x_{\min}$ that minimises $D$.
  \item Discard fits with fewer than 15 tail points.
\end{enumerate}

If the \texttt{powerlaw} package is installed it is used instead of the
built-in fallback (same methodology; the package is the canonical
Clauset~2009 reference implementation).

\subsection{Power-law versus log-normal likelihood-ratio test (LRT)}
\label{sec:lrt}

For each aggregated (geometry, damp\_idx) sample, a Vuong-style
likelihood-ratio test distinguishes whether the data favours a power law
over a log-normal, following Clauset~2009 \S4:
\[
  R = \sum_{i} \ln\frac{p_{\mathrm{PL}}(s_i)}{p_{\mathrm{LN}}(s_i)}, \qquad
  p = \operatorname{erfc}\!\left(\frac{|R|}{\sqrt{2}\,\hat{\sigma}_{R}\sqrt{n}}\right)
\]
$R > 0$ and $p < 0.10$ is recorded as ``power-law preferred''.  The columns
\texttt{lrt\_R}, \texttt{lrt\_p}, and \texttt{lrt\_favors\_powerlaw} appear in
\file{powerlaw\_fits.csv} and \file{powerlaw\_fits\_S.csv}.

\subsection{Waiting-time distribution between avalanches}
\label{sec:waiting_time}

From the S-column inter-arrival times, the analysis fits two models:

\begin{description}
  \item[Exponential (Poisson process).]  MLE: $\hat{\lambda} = 1/\langle\tau\rangle$.
  \item[Power law.]  MLE as above.
\end{description}

The reported metrics include:
\begin{itemize}
  \item Mean inter-arrival time $\langle\tau_w\rangle$.
  \item Coefficient of variation $\mathrm{CV} = \sigma_{\tau_w}/\langle\tau_w\rangle$
    ($\mathrm{CV} \gg 1$ signals bursty / SOC-like timing).
  \item LRT statistic: $\Lambda = 2(\mathcal{L}_{\mathrm{PL}} - \mathcal{L}_{\mathrm{Exp}})$;
    positive values indicate departure from Poisson.
\end{itemize}

\subsection{Temporal autocorrelation of $M(t)$}
\label{sec:autocorr}

Motivated by the 2025 echo-state / memory-capacity result for ASI
(\textit{Taniguchi, Sci. Rep.\ 2025}), the script computes the
normalised autocorrelation function
\begin{equation}
  C(\ell) = \frac{\langle M_c(t)\,M_c(t+\ell)\rangle}{\mathrm{var}(M)}
\end{equation}
for lags up to $25\%$ of the run length.  The correlation time is
\begin{equation}
  \tau_{\mathrm{corr}} = \int_0^{\ell^*} C(\ell)\,d\ell
\end{equation}
where $\ell^*$ is the lag of the first zero-crossing of $C$.
The prediction is that $\tau_{\mathrm{corr}}$ peaks near $Q \approx 1$
(critical damping).

\subsection{Output files}

\begin{longtable}{p{5.5cm}p{9.5cm}}
\toprule
File & Contents \\
\midrule
\endhead
\file{summary\_hysteresis.csv} & Per-run metrics: area, $B_c$, $M_r$,
  avalanche counts (both methods), waiting-time statistics, $\tau_{\mathrm{corr}}$ \\
\file{summary\_relax.csv} & Per-run domain statistics from Stage~2 \\
\file{aggregated\_loop\_metrics.csv} & Mean $\pm$ std of area, $B_c$, $M_r$
  aggregated over seeds \\
\file{powerlaw\_fits.csv} & MLE results ($\tau$, $\tau_{\mathrm{err}}$, $x_{\min}$, KS,
  LRT) using the $dM/dB$ avalanche sizes \\
\file{powerlaw\_fits\_S.csv} & Same, using the spin-flip counter (S-column) sizes \\
\file{area\_vs\_Q.png} & Loop area vs.\ $Q$ (by geometry) \\
\file{Bc\_vs\_Q.png} & Coercive field vs.\ $Q$ \\
\file{Mr\_vs\_Q.png} & Remanence vs.\ $Q$ \\
\file{tau\_vs\_Q.png} & \textbf{Central result:} avalanche exponent $\tau$ vs.\ $Q$
  ($dM/dB$ method) \\
\file{tau\_vs\_Q\_Scol.png} & Same, using S-column sizes (cross-check) \\
\file{avalanche\_histograms.png} & Log-log $P(s)$ histograms with MLE fit lines \\
\file{domain\_size\_vs\_Q.png} & Mean domain size and domain count vs.\ $Q$ \\
\file{waiting\_time\_vs\_Q.png} & Mean IAT and LRT statistic vs.\ $Q$ \\
\file{autocorr\_vs\_Q.png} & Correlation time $\tau_{\mathrm{corr}}$ of $M(t)$ vs.\ $Q$ \\
\bottomrule
\end{longtable}

\subsection{Analysis command-line interface}

\begin{tabular}{p{4cm}p{2.2cm}p{8.5cm}}
\toprule
Flag & Default & Description \\
\midrule
\param{sweep\_dir} & \textit{required} & Directory written by a sweep script \\
\param{out\_dir} & \textit{required} & Where to save analysis results \\
\param{mad\_threshold} & 4.0 & MAD multiplier $\kappa$ for Method~1 avalanche detection \\
\param{forc\_dir} & \textit{None} & (Stub) Directory for FORC output; not yet implemented \\
\bottomrule
\end{tabular}

% ======================================================================
\section{Practical Recipes}
\label{sec:recipes}
% ======================================================================

\subsection{Pipeline verification (quick test)}

\begin{lstlisting}
# Serial (V01): 12x12 grid, 1 seed, 3 dampings, 8 periods -- finishes in seconds
python3 damping_sweep.py --out_dir "$(pwd)/quicktest" --quick_test 1

# Parallel (V02, 2 workers): same reduced parameters
rm -rf /tmp/bartest
python3 damping_sweepV02.py \
    --out_dir /tmp/bartest --quick_test 1 \
    --n_workers 2 --geometries square

# Analyse immediately
python3 damping_sweep_analysis.py \
    --sweep_dir /tmp/bartest \
    --out_dir   /tmp/bartest/analysis
\end{lstlisting}

\noindent During the V02 quick test, six jobs complete and you should see the
progress bar smoothly fill from~0\% to~100\% on a single terminal line, with
the active-worker countdown decreasing in the suffix.

\subsection{Standard campaign (CPU, parallel)}

\begin{lstlisting}
# 8 dampings x 3 geometries x 5 seeds = 240 hysteresis runs + 240 relax
python3 damping_sweepV02.py \
    --out_dir "$(pwd)/results" \
    --n_workers 4

# Post-process
python3 damping_sweep_analysis.py \
    --sweep_dir "$(pwd)/results" \
    --out_dir  "$(pwd)/results/analysis"
\end{lstlisting}

\subsection{Resume an interrupted campaign}

\begin{lstlisting}
python3 damping_sweepV02.py \
    --out_dir "$(pwd)/results" \
    --n_workers 4 --resume 1
\end{lstlisting}

\subsection{GPU machine}
\label{sec:gpu}

On a machine with a CUDA GPU and CuPy installed, use \param{use\_gpu~1}
with \textbf{exactly one worker}.  The GPU accelerates the dipolar
tensor matrix--vector product inside \code{relax()} (the ``step
matriz-vetor~V40'' operation), typically by a factor of 10--50$\times$
for a 30$\times$30 array; one fast serial worker beats multiple slow
CPU workers.

\warn{Never set both \texttt{--use\_gpu~1} and \texttt{--n\_workers~$>$~1}:}
multiple processes cannot share a single CUDA context and the campaign
will crash.

\begin{lstlisting}
# Verify CuPy installation first
python3 -c "import cupy; print('CuPy OK, devices:', cupy.cuda.runtime.getDeviceCount())"

# Quick GPU test (6 jobs, should complete in a few seconds on any GPU)
rm -rf /tmp/gpu_test
python3 damping_sweepV02.py \
    --out_dir /tmp/gpu_test --quick_test 1 \
    --geometries square --use_gpu 1 --n_workers 1

# Full production campaign on GPU
python3 damping_sweepV02.py \
    --out_dir ~/results/law3m_gpu \
    --grid_n 30 --n_seeds 8 --n_dampings 14 \
    --Q_min 0.02 --Q_max 30.0 \
    --geometries square,triangular,honeycomb \
    --t_sim_periods 80.0 --B_max_factor 8.0 \
    --damping_noise 0.05 --dt_factor 0.04 \
    --use_gpu 1 --n_workers 1 --resume 1
\end{lstlisting}

\noindent Install CuPy matching your CUDA version if not yet present:
\begin{lstlisting}
pip install cupy-cuda12x   # or cupy-cuda11x, cupy-cuda118, etc.
\end{lstlisting}

\subsection{LAW3M 2026 --- Recommended production campaign}
\label{sec:law3m}

This is the parameter set designed to produce the key figures for the
LAW3M~2026 conference contribution on ``Inertial effects on avalanche
universality classes in classical dipolar lattices''.

\begin{lstlisting}
# STAGE 1: Run the campaign overnight (approx. 6-8 h on 4 cores)
# 14 Q-values x 3 geometries x 8 seeds = 336 hysteresis + 336 relax runs

python3 damping_sweepV02.py \
    --out_dir  ~/results/law3m_2026 \
    --grid_n   30            \  # 30x30 = 900 needles per run
    --n_seeds  8             \  # 8 noise realisations per (geometry,Q)
    --n_dampings 14          \  # 14 Q-points from Q=0.02 to Q=30
    --Q_min    0.02          \  # deep overdamped: monotonic M(B)
    --Q_max    30.0          \  # deep underdamped: inertial overshoot
    --geometries square,triangular,honeycomb \
    --t_sim_periods 80.0     \  # 80 T0 per cycle (5 segments)
    --B_max_factor  8.0      \  # B_max = 8 x B_ref (ensures saturation)
    --damping_noise 0.05     \  # 5% per-needle variation (realistic disorder)
    --dt_factor 0.04         \  # finer step; important at large Q
    --n_workers 4            \
    --resume 1

# STAGE 2: Analysis
python3 damping_sweep_analysis.py \
    --sweep_dir ~/results/law3m_2026 \
    --out_dir   ~/results/law3m_2026/analysis \
    --mad_threshold 3.5
\end{lstlisting}

\medskip
\noindent\textbf{Key figures for the conference presentation:}
\begin{enumerate}
  \item \file{tau\_vs\_Q.png} and \file{tau\_vs\_Q\_Scol.png} ---
    the central physics result: does $\tau$ vary with $Q$? Does it differ
    by geometry?
  \item \file{autocorr\_vs\_Q.png} --- memory/correlation time vs.\ $Q$;
    tests the echo-state prediction.
  \item \file{avalanche\_histograms.png} --- visual $P(s)$ comparison between
    the underdamped ($Q \approx 30$) and overdamped ($Q \approx 0.02$) limits.
  \item Side-by-side video frames from \texttt{compass.py} direct calls (see
    below) illustrating chain-reaction vs.\ sluggish switching.
\end{enumerate}

\subsection{Demo animations (underdamped vs.\ overdamped)}

\begin{lstlisting}
# Underdamped (Q~15): inertial overshoot, chain-reaction avalanches
python3 compass.py --N 20 --M 20 --geometry square \
    --field_mode hysteresis --B_ext 0.003 --t_sim 25 \
    --damping 2.4e-7 --dt_factor 0.04 \
    --frame_dir ~/frames/underdamped --frame_every 5

# Overdamped (Q~0.05): sluggish monotonic switching
python3 compass.py --N 20 --M 20 --geometry square \
    --field_mode hysteresis --B_ext 0.003 --t_sim 25 \
    --damping 7.0e-5 --dt_factor 0.04 \
    --frame_dir ~/frames/overdamped --frame_every 5
\end{lstlisting}

% ======================================================================
\section{Physical Parameter Reference}
\label{sec:params}
% ======================================================================

Default values for a tabletop steel needle with $R = 0.025$\,m and
$f = 0.80$:

\begin{center}
\begin{tabular}{llll}
\toprule
Quantity & Symbol & Value & Unit \\
\midrule
Needle half-length & $L/2$ & 20.0 & mm \\
Needle width & $W$ & 4.4 & mm \\
Needle thickness & $e$ & 0.4 & mm \\
Steel density & $\rho$ & 7850 & kg\,m$^{-3}$ \\
Saturation magnetisation & $M_s$ & $1.59 \times 10^6$ & A\,m$^{-1}$ \\
Magnetic moment (derived) & $m$ & $\approx 5.6 \times 10^{-5}$ & A\,m$^2$ \\
Moment of inertia (derived) & $I$ & $\approx 7.6 \times 10^{-9}$ & kg\,m$^2$ \\
Nearest-neighbour distance & $r_{\mathrm{nn}}$ & 50 & mm \\
Dipolar reference field & $B_{\mathrm{ref}}$ & $\approx 0.36$ & mT \\
Natural frequency & $\omega_0$ & $\approx 22.8$ & rad\,s$^{-1}$ \\
Natural period & $T_0$ & $\approx 0.276$ & s \\
Default damping & $b$ & $5\times 10^{-8}$ & N\,m\,s\,rad$^{-1}$ \\
Corresponding $Q$ & & $\approx 3.5$ & \\
\bottomrule
\end{tabular}
\end{center}

% ======================================================================
\section{Notes and Caveats}
\label{sec:caveats}
% ======================================================================

\begin{itemize}
  \item \textbf{Avalanche statistics require large samples.}  A single 30$\times$30
    hysteresis run may yield only $\mathcal{O}(10)$--$\mathcal{O}(100)$ events.
    The analysis aggregates over all seeds before the MLE fit; fewer than
    8 seeds at any (geometry, $Q$) point typically gives unstable $\tau$ estimates.

  \item \textbf{Method~2 (S-column) preferred.}  The $dM/dB$ MAD method is sensitive
    to the threshold multiplier $\kappa$ and to numerical noise near field
    inflection points.  The S-column method counts discrete needle reversals
    with no free parameters; use \file{tau\_vs\_Q\_Scol.png} as the primary result
    and \file{tau\_vs\_Q.png} as a cross-check.

  \item \textbf{Do not mix GPU with $n_{\mathrm{workers}} > 1$.}  Multiple worker
    processes cannot safely share a CUDA context.  See
    Section~\ref{sec:gpu} for the correct GPU workflow.

  \item \textbf{Progress bar requires a real terminal.}  The \texttt{\textbackslash r}
    overwrite trick works only when stdout is connected to an interactive
    terminal (TTY).  When output is piped (e.g.\ \texttt{| tee log.txt} or
    redirected to a file), each bar update appears on a new line in the file.
    This is cosmetic only; the campaign data are unaffected.

  \item \textbf{Pivot parameters default to zero.}  If the physical experiment uses
    a metallic pivot pin that contributes to $I$, specify \param{pivot\_radius},
    \param{pivot\_thickness}, and \param{pivot\_density} (or \param{pivot\_mass})
    to maintain consistency between the simulation and the tabletop apparatus.

  \item \textbf{The cutoff is in units of $r_{\mathrm{nn}}$.}  The \param{cutoff}
    argument passed to \code{relax()} is expressed as a multiple of the
    nearest-neighbour distance ($r_{\mathrm{nn}} = 2R$), \textit{not} in metres.
    The default value of 3.5 includes up to $\approx$third-shell neighbours.
\end{itemize}

\end{document}
